\documentclass[11pt]{exam}
\usepackage[left=1in,right=0.75in,top=1in]{geometry}                		% See geometry.pdf to learn the layout options. There are lots.

\pagestyle{headandfoot}
\firstpageheader{\bfseries{Biogeography, Exam 2, F13}}{}{\bfseries{Name: \enspace \makebox[2.5in]{\hrulefill}}}
\runningheader{}{}{\small{pg. \thepage}}
%\runningheadrule

\footer{}{}{}%{\makebox[0.75in]{\hrulefill}}
\extrafootheight{-0.5in}

\pointsinmargin
\marginpointname{ pts}

\usepackage{multicol}
\usepackage{xcolor,graphicx}
	\graphicspath{{img/}} % set of paths to search for images

%% Remove comment %% to print answer key.
%\printanswers
\unframedsolutions
\renewcommand{\solutiontitle}{}
\SolutionEmphasis{\bfseries}

\begin{document}

\begin{questions}

\question[3]
About when (from - to) did the Pleistocene occur?

\question[5]
Explain the difference between obliquity and eccentricity.


\question[10]
Time to get serious: Describe the difference between the target effect and the rescue effect.

	\begin{solution}
	Albedo is amount of light reflected away from earth. More glaciers would have reflected back more sun, leading to cooler temperatures.
	\end{solution}

\vspace{3in}


\question[10]
Describe the albedo effect and how it contributed to expansion of Pleistocene glaciers.

	\begin{solution}
	Albedo is amount of light reflected away from earth. More glaciers would have reflected back more sun, leading to cooler temperatures.
	\end{solution}


%\question[10]
%Average summer temperatures during the Pleistocene have been linked to glacial extent (amount of glacial coverage) in the northern hemisphere. Explain how glacial extent might have changed if eccentricity decreased and obliquity  decreased at some point during the Pleistocene.  Describe how the coverage would have changed (more or less) and how these two cycles contributed to that change.

%	\begin{solution}
%	Decreasing eccentricity would bring the earth closer to the sun during the summer, causing warmer summers. Increasing obliquity would have further increased the amount of sunlight reaching the northern hemisphere. These together would have contributed to warmer global temperatures, decreasing the extent of glacial coverage.
%	\end{solution}


\newpage

\question[15]
Australia was not covered by glaciers but climate change during the Pleistocene still caused range contraction and expansion for many taxa.  An Australian gecko, \textit{Heteronotia binoei}, has the distribution indicated by the black dots in the inset below.  The large gap in the distribution suggests one of two scenarios: post-Pleistocene range expansion or a long-term barrier to dispersal between the two regions. Strasburg et al. (2007) tested these alternate hypotheses. Below is a simplified version of their parsimony network.  Based on the network, state which hypothesis is supported and describe how the network supports the hypothesis you chose. 

Each circle represents one allele. Color corresponds to the population sampled; grey is the western cluster and white is the central cluster. Each line connecting two circles indicates one mutation difference between alleles.  Small dashes indicate additional mutational differences between alleles.

	{\centering
	\includegraphics[width=\textwidth]{gecko_network}\par}

	\begin{solution}
	The western (grey) clade shows considerable genetic diversity, suggesting it is older and was likely a refuge. The central (white) clade has little genetic diversity, suggesting it is much younger and therefore the result of recent range expansion.
	\end{solution}

\newpage


\question[15]
Use Lomolino's extension of the MacArthur-Wilson model of island biogeography to explain why endemism is high on large, isolated islands, as shown in the figure below.  

\begin{minipage}{0.3\textwidth}
	\hfill
	\begin{solution}
	Large islands provide resources to support abundance, different niches for adaptation. Isolation means little to no gene flow preventing evolution.
	\end{solution}
\end{minipage}
\begin{minipage}{0.6\textwidth}
	\hfill \includegraphics[width=\textwidth]{lomolino_endemism}
\end{minipage}



\newpage

\question[20]
Abbott's Duiker is a small antelope endemic to the fragmented Eastern Arc forests of Tanzania.  Known localities with the duiker are shown in the inset.  Assume you are given the task of increasing survivorship and movement of individuals between the southernmost and northernmost locations.  Apply whatever island biogeography concepts you think are necessary to accomplish this goal.  You must explain \textit{how} each concept you choose helps achieve the goal.  You must also discuss how corridors can contribute, but you must also discuss at least one potential disadvantage of corridors (as discussed in class on Tuesday). I have extra paper if you need more writing space.

	{\centering
	\includegraphics[width=0.8\textwidth]{abbott_duiker_distrib}\par}

	\begin{solution}[1in]
	As the basin spreads, the taxa become less and less similar.
	\end{solution}



\end{questions}

\end{document}